\begin{statement}{C.2.UN.35}{Proved}{Every prescribed cyclic smooth-multiple elliptic quotient occurs in a zero-Euler torus filling}
Let $q:E\to B$ be a connected compact smooth elliptic
surface with a fixed elliptic Jacobian $F$, identity
integral $H_1$ monodromy, and only smooth multiple
fibres. Suppose $g(B)\le1$. Let $Q$ be a finite
cyclic group acting holomorphically on $E$, preserving
$q$, and acting faithfully on $B$. The trivial group
is allowed. The smooth minimal model of $E/Q$ occurs
as the minimal normalization model of a component of
a fibre in a smooth compact threefold over $\PP^1$
whose general fibre is a complex two-torus. It also
occurs in such a fibre with reduced-support Euler
number zero. The construction retains the given
surface and its given action before taking the
quotient: the distinguished component of the
constructed covering threefold is $Q$-equivariantly
isomorphic to $E$ over $B$.

Consequently the minimal normalization models of
nonnegative Kodaira dimension occurring in arbitrary
two-torus degenerations are exactly the following
two explicitly specified classes:
\begin{itemize}
\item the minimal resolutions of finite cyclic
quotients of complex two-tori that have nonnegative
Kodaira dimension;
\item the minimal models of $E/Q$ in the preceding
construction that have nonnegative Kodaira dimension.
\end{itemize}
The same two classes occur among zero-Euler fillings.
The torus-quotient class is described completely in
UN23 and UN32. In the second class, all multiplicities,
translation gluing, and the cyclic action are part
of the specified input; none is discarded by the
occurrence construction.
\end{statement}

\begin{proof}
We give the extension on actual finite root charts,
and then prove the exhaustive assertion.

\par\medskip\noindent\textbf{The actual orbicurve torsor.}
Write the multiple fibres as $m_iC_i$ over $p_i\in B$,
with $m_i>1$. The normalized root chart proved in
UN29 and (C.KX.3)--(C.KX.5) is
\[
 (F\times\Delta_u)/\langle\delta_i\rangle,
 \qquad
 \delta_i(z,u)=(z+\tau_i,\zeta_{m_i}u),
 \qquad \operatorname{ord}(\tau_i)=m_i,
 \qquad t=u^{m_i}.
                                                               \tag{C.EO.1}
\]
The free translation on $F$ is essential here.
Let $\mathcal B=B[(p_i,m_i)]$ denote the root
orbicurve, whose chart at $p_i$ is
$[\Delta_u/\mu_{m_i}]$ with $t=u^{m_i}$.
Changes of local origins of the fixed elliptic
curve are translations: the linear part is the
identity because the integral monodromy is the
identity. Thus these charts and the ordinary
product charts exhibit the given $E$ as a
principal $F$-torsor over $\mathcal B$.
Its total space is the original smooth surface,
because all the root inertia actions in (C.EO.1)
are free. The action of $Q$ induces the functorial
action on the root orbicurve, and its linear action
on $F$ is a constant finite character
\[
                 \beta:Q\longrightarrow
                            \operatorname{Aut}(F,0).
                                                               \tag{C.EO.2}
\]
The constancy follows from the discreteness of
$\operatorname{Aut}(F,0)$ and connectedness of
the regular base. The actual action on the
torsor includes its translation part.

\par\medskip\noindent\textbf{A lift that retains the whole action.}
Put $F=\C/\Lambda$ and let
$\ell:\Lambda\otimes_{\Z}\C\to\C$
be the complex-linear extension of the inclusion
of the period lattice. There is an exact sequence
of complex Lie groups
\[
 0\longrightarrow K\longrightarrow
 T=(\Lambda\otimes_{\Z}\C)/\Lambda
 \xrightarrow{\,\Phi\,}F\longrightarrow0,
 \qquad K=\ker\ell\simeq(\C,+),
 \qquad T\simeq(\C^*)^2.                              \tag{C.EO.3}
\]
For a basis $(\lambda_1,\lambda_2)$ of $\Lambda$,
the last map is explicitly
\[
 \Phi(z_1,z_2)=
 \frac{\lambda_1\log z_1+\lambda_2\log z_2}{2\pi i}
                    \pmod\Lambda .                 \tag{C.EO.4}
\]
Changing either logarithm changes the numerator
by the corresponding period. The map is locally
split and surjective. Its kernel is exactly
the image of the complex line $\ker\ell$:
a vector mapping to a period can first be
changed by that integral lattice vector.
The intersection $\ker\ell\cap\Lambda$ is zero,
and the kernel is closed because $F$ is Hausdorff.
This proves the exactness, including the actual
additive kernel, without a choice of a cyclic
subgroup of the elliptic periods.

Every $\beta(g)$ acts on $\Lambda$ by an integral
matrix and hence on $T$ by the corresponding
monomial automorphism. It preserves $K$; on
that line its character is $\beta(g)^{-1}$,
because its determinant on $\Lambda$ is one.
Sequence (C.EO.3) is therefore $Q$-equivariant.
On the quotient orbicurve $\mathcal C=[\mathcal B/Q]$
it gives an exact sequence of sheaves of
holomorphic sections
\[
 0\longrightarrow\mathcal K\longrightarrow
 \mathcal T\longrightarrow\mathcal F\longrightarrow0.
                                                               \tag{C.EO.5}
\]
Here $\mathcal K$ is a holomorphic line bundle
with its additive group law; the other sheaves
retain the indicated finite monodromy. The given
equivariant elliptic torsor is an element of
$H^1(\mathcal C,\mathcal F)$.

One has $H^2(\mathcal C,\mathcal K)=0$.
Indeed finite-group invariants on complex vector
spaces are exact by averaging. On each finite
root chart the higher invariant direct images
therefore vanish, and the invariant direct image
of a coherent sheaf is coherent on the coarse
curve. Coherent cohomology in degree two on a
compact complex curve is zero. The same averaging
for $Q$ proves the asserted vanishing on
$\mathcal C$. The long exact sequence of
(C.EO.5) consequently supplies an actual
$\mathcal T$-torsor $P$ and an isomorphism
\[
                  P\mathbin{\times^{T}}F
                           \simeq E                 \tag{C.EO.6}
\]
in the equivariant orbicurve category. In
particular the descent isomorphisms of $P$
satisfy the full cocycle identity; they are
not merely isomorphisms of underlying line
bundle classes.

\par\medskip\noindent\textbf{An invariant meromorphic frame.}
The faithful action on $B$ makes
$L=\C(B)$ a cyclic extension of
$k=\C(B/Q)$. At the generic point the root
inertia is absent. The monomial descent of
$(\C^*)^2_L$ defines an algebraic $k$-torus
$T_k$, and the generic fibre of $P$ is an
algebraic torsor under it. To see algebraicity
directly, pull back to $\mathcal B$ and choose
meromorphic sections of its two associated
orbifold line bundles. Such sections exist:
after separating the finitely many root weights,
this is the ordinary meromorphic-section theorem
for line bundles on a compact curve. The
equivariant descent isomorphisms are then
monomials multiplied by meromorphic functions
on $B$, which belong to $L$. Their existing
cocycle identities give precisely this algebraic
generic torsor.

It has a $k$-rational point. This is the
rank-two torus case of Steinberg's theorem
\cite[Theorem~1.9]{Steinberg1965}; the independent
curve-field and cyclic-lattice proof below
establishes the required vanishing for exactly
the representations used here.
\input{sections/equivariant_torus_h1.inc.tex}
The curve-field proof in that calculation also
checks the precise field hypothesis in the
printed statement of Steinberg's theorem.
It applies to every finite extension of $k$,
since that extension is another complex curve
function field. A central division algebra
of degree $d>1$ over any such field has a
reduced norm of degree $d$ in $d^2$ variables.
The proved $C_1$ property would give a nonzero
element of reduced norm zero, impossible in
a division algebra. Thus all the indicated
division algebras are commutative, which is
the cohomological-dimension-at-most-one
condition used in \cite[Theorem~1.9]{Steinberg1965}.
Pull back that rational point to $L$.
It is a meromorphic section $s$ of $P$ on
$\mathcal B$ satisfying the exact equality
\[
                         g(s(b))=s(gb)
                              \qquad(g\in Q).        \tag{C.EO.7}
\]
Thus the translation cocycle has been retained
in a genuine invariant meromorphic frame.

Let $M=\operatorname{Hom}(T,\C^*)$ be the
character lattice. For $\chi\in M$, write
$L_\chi$ for the associated line bundle and
$s_\chi$ for its meromorphic section induced
by $s$. Its divisor defines an additive map
\[
 D:M\longrightarrow\operatorname{Div}(\mathcal B),
       \qquad D(\chi)=\operatorname{div}(s_\chi).
                                                               \tag{C.EO.8}
\]
This map is $Q$-equivariant with its actual
action on characters and divisors, by (C.EO.7).
Choose two basis characters. Their supports,
together with the root points, lie in one
finite $Q$-invariant subset $\Sigma\subset B$.
In terms of coarse points, every coefficient
of $D(\chi)$ is integral away from the root
points and belongs to $m_i^{-1}\Z$ at $p_i$.
We have the canonical multiplicative
identifications
\[
          (L_\chi,s_\chi)
                   \simeq(\OO_{\mathcal B}(D(\chi)),1),
                                                               \tag{C.EO.9}
\]
where $1$ denotes the canonical meromorphic
divisor section. They commute with $Q$ and
with addition of characters. There is no
remaining scalar linearization or group
two-cocycle in (C.EO.9).

\par\medskip\noindent\textbf{An equivariant vertical extension.}
We construct a smooth projective surface
$V_0\to\PP^1_t$ containing $B$ in its
central fibre and extending its given
$Q$-action. Let $n=|Q|$.
For elliptic $B$, choose any faithful
character $\xi:Q\to\C^*$ when $n>1$
(and $\xi=1$ for the trivial group),
and take $V_0=B\times\PP^1_t$
with action $(b,t)\mapsto(gb,\xi(g)t)$,
and use $B\times\{0\}$.
For $B\simeq\PP^1$, first diagonalize
its given cyclic action as
$z\mapsto\lambda(g)z$ and define
$\xi=\lambda^{-1}$. Faithfulness makes
this a faithful disc character too.
For the trivial group use both
characters equal to one. Choose any
elliptic curve $A$ and point $o\in A$.
Blow up $(o,0)$ in $A\times\PP^1_t$,
where the action is
$(a,t)\mapsto(a,\xi(g)t)$.
The exceptional curve is identified with
the given $B$: in a local coordinate $x$
at $o$, its affine coordinate $z=x/t$
has exactly the required action. The
other chart supplies its point at infinity.

Blow up every point of $\Sigma$ on the
distinguished strict curve $B$, in complete
$Q$-orbits. Call the resulting disjoint
exceptional curves $T_p$, $p\in\Sigma$,
and call the new surface $V$. The strict
curve is still isomorphic to $B$, and
\[
 \OO_V(T_p)|_B=\OO_B(p),\qquad
                     g(T_p)=T_{gp}.                \tag{C.EO.10}
\]
These identities include the canonical
sections and their equivariant structures.
There is no restriction at a fixed point
or at the node of the original central
fibre. At a node with coordinates $xy=0$
and $B=(x=0)$, its blowup has chart
$x=uv,y=v$. Here the strict $B$ is
$u=0$, $T_p$ is $v=0$, and the
restriction of its canonical section
is the simple parameter $v$ on $B$.
This proves (C.EO.10) also in that case.
The other central branch meets $T_p$
at the other point of this exceptional
curve; it does not meet the strict $B$.

Root $V$ along each $T_{p_i}$ to order
$m_i$, obtaining the orbifold surface
$\mathcal V$. Its local chart at a
point of $T_{p_i}$ is
\[
 [\Delta_u\times\Delta_v/\mu_{m_i}],
           \qquad x=u^{m_i},\qquad
           \zeta\cdot(u,v)=(\zeta u,v).
                                                               \tag{C.EO.11}
\]
The $T_{p_i}$ are disjoint, so there
is only this one inertia group in
each such chart. Its restriction to
$B$ is precisely the previous root
orbicurve $\mathcal B$, including the
functorial $Q$-action. The transverse
intersection in (C.EO.10) is what
identifies its root line and root
section with those on $\mathcal B$.

Replace each coarse point $p$ in
$D(\chi)$ by $T_p$, with the same
coefficient. Denote the resulting
orbifold Cartier divisor by
$\widetilde D(\chi)$. All its
coefficients are legitimate integral
root-divisor coefficients on
$\mathcal V$. The map
$\chi\mapsto\widetilde D(\chi)$ is
additive and $Q$-equivariant because
the replacement $p\mapsto T_p$ is.
The character lines
$\OO_{\mathcal V}(\widetilde D(\chi))$
therefore define a principal $T$-torsor
$\widetilde P$ with its exact monomial
$Q$-action. Define the principal
elliptic bundle
\[
              \widetilde Y
                =\widetilde P\mathbin{\times^T}F
                     \longrightarrow\mathcal V .   \tag{C.EO.12}
\]
Concretely, on every line-bundle
trivialization its transition map is
translation on $F$ by the value under
(C.EO.4) of the two multiplicative
transition functions. Equations
(C.EO.9)--(C.EO.10) give an actual
$Q$-equivariant isomorphism
\[
                   \widetilde Y|_{\mathcal B}
                                   \simeq E,         \tag{C.EO.13}
\]
over the given base, not merely
an equality of topological Chern
classes. Locally the meromorphic
frame identification extends
holomorphically because the two
divisors agree: their quotient is
a nowhere-zero holomorphic unit.
The same assertion holds in the
root coordinates. Since both
frames in (C.EO.9) are exactly
equivariant, these local isomorphisms
also retain the given $Q$-action.
More explicitly, a lift of a base
map to $x=u^{m_i}$ involves a
choice of an $m_i$-th root of
its nonzero local multiplier.
Two choices differ by the
root inertia in (C.EO.11).
The character lines simultaneously
apply the corresponding elliptic
translation in (C.EO.14).
Hence the two lifts induce the
same map on the total coarse
quotient. Compositions of lifts
agree there for the same reason.
This is the finite-chart form
of the canonical section-preserving
action, and proves that no local
root choice changes its group law.

\par\medskip\noindent\textbf{Smoothness, properness and the generic torus.}
Let $(a_1,a_2)\pmod {m_i}$ be the
two weights of $P$ at $p_i$.
In (C.EO.11) the total elliptic
bundle has the finite chart
\[
 (\Delta_u\times\Delta_v\times F)/\mu_{m_i},
 \quad
 \zeta_{m_i}(u,v,z)
     =(\zeta_{m_i}u,v,z+\tau_i),\quad
 \tau_i=\frac{a_1\lambda_1+a_2\lambda_2}{m_i}
                                           \pmod\Lambda .
                                                               \tag{C.EO.14}
\]
The order of $\tau_i$ is exactly
$m_i$ by (C.EO.1) and (C.EO.6).
Every nonidentity inertia element
therefore acts freely on the total
chart, including $u=0$ and every
central crossing $v=0$. Thus the
coarse space in (C.EO.12) is a
smooth compact complex threefold;
we continue to call it
$\widetilde Y$. Properness follows
on these finite charts from
properness of the compact factor
$F$ and glues over the compact
surface $V$. The principal-bundle
transition functions in (C.EO.12)
give the gluing maps explicitly.

All divisors used to construct
$\widetilde P$ are vertical over
$t=0$. Their canonical meromorphic
sections are therefore nowhere
zero on the complement of that
fibre. There (C.EO.12) is the
actual product
\[
 \widetilde Y|_{t\ne0}
       \simeq A\times F\times(\PP^1\setminus\{0\}),
                                                               \tag{C.EO.15}
\]
where $A=B$ in the elliptic-base
case and $A$ is the chosen auxiliary
elliptic curve in the rational-base
case. The action is the prescribed
action on $A$ (the identity in the
rational-base construction),
the linear action $\beta$ on $F$,
and multiplication by $\xi$ on
$t$. This formula follows from
the invariant divisor frames,
so no additional translation
cocycle survives off zero.

Take the finite quotient by $Q$
and resolve its singularities.
Finite holomorphic actions can
be linearized locally, and these
cyclic quotients have compatible
proper resolutions; a resolution
of the compact quotient, chosen
over its singular locus, is
enough. Its map to
$\PP^1_T$, $T=t^n$, has general
fibre exactly $A\times F$:
the $n$ covering fibres are
permuted simply transitively
for $T\ne0,\infty$. The quotient
is smooth there, so the
resolution leaves those fibres
unchanged. The Stein factor is
a finite map of degree one
to the normal curve $\PP^1$,
and is the identity. Thus
the resulting proper map has
connected fibres.

The distinguished component
before the quotient is exactly
the given $E$ by (C.EO.13).
Its image in the finite quotient
is $E/Q$: taking invariants
is exact in characteristic zero,
so the restriction of the
quotient structure sheaf to
this image is
$(\OO_E)^Q$. It is normal
because $E$ is normal. The
strict component on the ambient
resolution maps properly and
birationally to $E/Q$.
Resolving that component and
contracting its surface
$(-1)$-curves gives exactly
the prescribed minimal model.
Finally UN1 gives even Euler
number for the reduced support.
The actual vertical point and
quartic blowups of UN4 change
that number by $+2$ and $-2$
respectively. Finitely many
such operations make it zero
and preserve the minimal
normalization model through
their explicit proper
birational maps. This proves
the occurrence assertion.

\par\medskip\noindent\textbf{Exhaustion of the nonnegative classes.}
Let $Z$ be a minimal normalization
model of nonnegative Kodaira
dimension in an arbitrary
two-torus degeneration. UN31
provides its actual cyclic
component cover and equivariant
minimal contraction. A torus
cover gives exactly the first
class in the statement. In
the remaining case UN22 and
UN29 give a fixed-Jacobian,
identity-monodromy fibration
with only smooth multiples
over a rational or elliptic
curve. We specify why an
appropriate such fibration
is preserved by the cyclic
action.

For a cover of Kodaira
dimension one it is the
intrinsic connected Iitaka
fibration. In algebraic
dimension one its base is
the intrinsic meromorphic
function field, and the
equivariant Stein factor
of UN20 applies. The
non-torus projective case
of Kodaira dimension zero
is bielliptic. Use its
actual finite free product
description
$(F_1\times F_2)/\Gamma$,
with $\Gamma$ acting by
translations on $F_1$
and with a nontrivial
linear action on $F_2$.
The invariant holomorphic
vector fields on the
product are exactly the
constant fields in the
$F_1$ direction. These
integrate to the connected
automorphism group of
the quotient, an elliptic
quotient of $F_1$.
Its orbit map is
\[
 (F_1\times F_2)/\Gamma
                   \longrightarrow F_2/\Gamma\simeq\PP^1.
                                                               \tag{C.EO.16}
\]
Every automorphism normalizes
the connected automorphism
group and hence preserves
this orbit fibration. Its
local stabilizers act by
free translations on
$F_1$; its root charts
are (C.EO.1). Therefore
it has the required
identity monodromy and
only smooth multiples.
This is the preserved
fibration used in the
bielliptic case.

Let $H$ be the kernel of
the cyclic action on the
chosen base. Its action
on a general elliptic
fibre is by translations:
a nontrivial linear
elliptic automorphism
has rational quotient,
which would make the
quotient surface ruled
and annihilate all
positive pluricanonical
sections, contrary to
$\kappa(Z)\ge0$.
The exact translation
and reflection charts
(C.EQ.3)--(C.EQ.4)
show that quotienting
the cover by $H$ is
again a smooth elliptic
surface with fixed
Jacobian, identity
monodromy, and only
smooth multiples.
The residual cyclic
group acts faithfully
on the base. Its
quotient is the same
actual surface quotient
as before, through
the finite quotient
maps. It is therefore
one of the inputs
already constructed
above.

Conversely, the
construction proves
occurrence for every
member of the second
class. For the first,
take the compact
diagonal quotient
$(A\times\PP^1_t)/G$
using a faithful
character on $t$,
resolve it, and use
the same UN4 balancing;
the exact component
comparison is UN24.
All occurrences may
therefore be chosen
with zero-Euler
support, proving
both asserted
equalities of classes.
\end{proof}


% C.2.UN.64 from part_C2_exceptional_leaves.tex

